% =============================================================
% Annexe A — Études comparatives détaillées
% Alimentée depuis Ch.1 (§1.2.3, §1.3.3) et Ch.2 (§2.3.4, §2.5, §2.6).
% Les grilles qui portent une justification de choix ont été remontées dans le corps
% du chapitre 2 ; cette annexe ne conserve que le détail de grille et les sources.
% =============================================================
\chapter{Études comparatives détaillées}
\label{ann:etudes-comparatives}

% --- Blancs autour des flottants, resserrés pour les annexes A à F
%     (valeurs LaTeX par défaut : 20 pt, 12 pt, 12 pt). Rétablis en fin de
%     l'annexe~F : le réglage ne sort pas du périmètre de cette équipe. ---
\setlength{\textfloatsep}{10pt plus 2pt minus 2pt}
\setlength{\intextsep}{8pt plus 2pt minus 2pt}
\setlength{\floatsep}{8pt plus 2pt minus 2pt}


\section{Audit de l'existant et positionnement (issus du \cref{ch:cadre-existant})}
\label{ann:etudes-comparatives-ch1}

Cette section porte le détail de ce que le \cref{ch:cadre-existant} condense : le
positionnement à quinze critères des trois modèles d'architecture comparés
(\cref{subsec:positionnement-modeles}). La répartition des 41~critères de l'audit du
29/07/2026 par catégorie et par statut n'y est pas redonnée : elle est portée par la
\cref{graph:audit-41} du corps, qui la publie en entier --- fonctionnel prêt à 86~\%, technique
et conformité à 0~\%, sécurité à mi-chemin, portant à la fois les points forts applicatifs et
les six constats bloquants.

\subsection{Positionnement complet à quinze critères}

Le \cref{tab:positionnement-15-complet} reprend l'intégralité de la comparaison introduite au
\cref{subsec:positionnement-modeles} entre le modèle A (architecture périmétrique de l'existant),
le modèle B (hébergement cloud managé sans démarche de sécurité explicite) et le modèle C (socle
Zero Trust visé par le projet). Les sept premiers critères proviennent des principes Zero
Trust~\cite{nist800207} ; les critères 8 à 11 reprennent les carences C2 à C4 et l'apport original
du projet ; les critères 12 à 15 traduisent des contraintes d'ingénierie réelles, sur lesquelles le
modèle C ne l'emporte pas systématiquement. Légende : \pictoPleinAudit~couvert par conception,
\pictoPartielAudit~partiellement couvert ou dépendant de la configuration,
\pictoVideAudit~non couvert.

\begingroup
\small
\setlength{\LTpre}{4pt}
\setlength{\LTpost}{4pt}
\begin{longtable}{|P{0.7cm}|P{7.6cm}|P{3.0cm}|c|c|c|}
\caption[Positionnement des trois modèles]{Positionnement complet des trois modèles sur quinze critères}
\label{tab:positionnement-15-complet}\\
\hline
\textbf{\#} & \textbf{Critère} & \textbf{Origine} & \textbf{A} & \textbf{B} & \textbf{C} \\
\hline
\endfirsthead
\hline
\textbf{\#} & \textbf{Critère} & \textbf{Origine} & \textbf{A} & \textbf{B} & \textbf{C} \\
\hline
\endhead
\hline
\endfoot
1 & Vérification explicite de l'identité à chaque requête & NIST SP 800-207 &
  \pictoVideAudit & \pictoPartielAudit & \pictoPleinAudit \\ \hline
2 & Aucune confiance accordée à la position réseau & NIST SP 800-207 &
  \pictoVideAudit & \pictoPartielAudit & \pictoPleinAudit \\ \hline
3 & Moindre privilège, une identité par charge de travail & NIST SP 800-207 &
  \pictoVideAudit & \pictoPartielAudit & \pictoPleinAudit \\ \hline
4 & Micro-segmentation, limitation des mouvements latéraux & NIST SP 800-207 &
  \pictoVideAudit & \pictoPartielAudit & \pictoPartielAudit \\ \hline
5 & Hypothèse de compromission : tout est journalisé & NIST SP 800-207 &
  \pictoVideAudit & \pictoPartielAudit & \pictoPleinAudit \\ \hline
6 & Collecte, corrélation et conservation des événements & NIST SP 800-207 &
  \pictoVideAudit & \pictoPartielAudit & \pictoPleinAudit \\ \hline
7 & Surveillance continue et amélioration & NIST SP 800-207 &
  \pictoVideAudit & \pictoVideAudit & \pictoPartielAudit \\ \hline
8 & Infrastructure entièrement décrite en code & Carence C2 &
  \pictoVideAudit & \pictoVideAudit & \pictoPleinAudit \\ \hline
9 & Contrôles de sécurité bloquants dans la livraison & Carence C3 &
  \pictoVideAudit & \pictoPartielAudit & \pictoPartielAudit \\ \hline
10 & Règles de détection écrites, versionnées et mesurées & Carence C4 &
  \pictoVideAudit & \pictoVideAudit & \pictoPartielAudit \\ \hline
11 & Priorisation des vulnérabilités par la menace observée & Apport du projet &
  \pictoVideAudit & \pictoVideAudit & \pictoPartielAudit \\ \hline
12 & Coût mensuel d'exploitation & Contrainte projet &
  \pictoPleinAudit & \pictoPartielAudit & \pictoPartielAudit \\ \hline
13 & Simplicité initiale, délai de mise en service & Contrainte projet &
  \pictoPleinAudit & \pictoPartielAudit & \pictoVideAudit \\ \hline
14 & Mise à l'échelle horizontale & Contrainte projet &
  \pictoVideAudit & \pictoPleinAudit & \pictoPartielAudit \\ \hline
15 & Maturité, support éditeur, compétences disponibles & Contrainte projet &
  \pictoPartielAudit & \pictoPleinAudit & \pictoVideAudit \\
\end{longtable}
\endgroup

\textbf{Lecture du tableau.} Les onze premières lignes montrent l'apport du modèle retenu, les
quatre dernières ce qu'il coûte, et se lisent avec la même attention. Le modèle~C perd sur la
simplicité initiale (ligne~13), sur la mise à l'échelle (ligne~14, hors objectif du projet) et sur
la maturité opérationnelle (ligne~15) ; le coût mensuel (ligne~12) n'est comparable qu'en
apparence, la facturation directe du modèle~A n'intégrant ni sauvegarde externalisée, ni détection,
ni reprise après incident. Ces trois critères perdus sont ce qui rend les onze autres crédibles, et
ils sont repris tels quels en \cref{sec:limites}.

\textbf{La ligne~4 appelle une note : son verdict n'a pas bougé, sa raison a changé du tout au
tout.} Avant la refonte du plan réseau, le modèle~C n'obtenait qu'un \pictoPartielAudit~faute de
toute segmentation ; celle-ci existe désormais (\cref{subsec:reseau-cible}) et le verdict reste
pourtant intermédiaire, pour un motif entièrement différent : \textbf{elle s'arrête au plan du
réseau}. Base de données et entrepôt de supervision demeurent partagés, le pare-feu autorisant
les deux identités à joindre la même adresse ; l'isolation du plan de la donnée est une
conception datée et chiffrée, non un état appliqué.

\section{État de l'art et choix technologiques (issus du \cref{ch:etat-art})}
\label{ann:etudes-comparatives-ch2}

Les grilles qui portent directement une justification de choix figurent dans le corps du
\cref{ch:etat-art} ; cette annexe ne conserve que ce qui alourdirait le corps sans y ajouter
d'argument.

\subsection{Comparaison des modes d'exécution}
\label{ann:modes-execution}

Trois modes d'exécution ont été comparés sur six critères. La conclusion étant déjà portée par le
\cref{tab:plateforme-10-criteres} et par la décision D01, le tableau en documente ici le seul
détail : l'exécution de conteneurs sans serveur réduit la charge d'administration et le coût à
faible trafic, introduit le démarrage à froid mesuré au \cref{ch:validation} et ne donne
qu'une isolation réseau intermédiaire.

Cette dernière ligne appelle une précision : son motif a changé sans que son verdict bouge. La
sortie réseau directe a levé la réserve initiale, chaque instance portant désormais une interface
que les règles nomment (\cref{subsec:reseau-cible}). Le pictogramme reste intermédiaire pour deux
motifs qui subsistent : le sélecteur disponible est l'étiquette réseau et non le compte de
service, et trois charges de l'application hébergée partagent un même compte, donc une même
étiquette, et restent indiscernables pour une règle --- un orchestrateur managé conserve
l'avantage sur ce seul critère.

\begin{table}[htbp]
\caption{Comparaison des trois modes d'exécution étudiés}
\centering
\scriptsize
{Légende : \pictoPleinAudit~favorable, \pictoPartielAudit~intermédiaire, \pictoVideAudit~défavorable.
\par\smallskip}
\begin{tabular}{|P{6.1cm}|P{3.0cm}|P{3.0cm}|P{3.0cm}|}
\hline
\textbf{Critère} & \textbf{Machines virtuelles} & \textbf{Orchestrateur managé} &
\textbf{Conteneurs sans serveur} \\
\hline
Charge d'administration & \pictoVideAudit & \pictoVideAudit & \pictoPleinAudit \\ \hline
Coût à faible trafic & \pictoPartielAudit & \pictoVideAudit & \pictoPleinAudit \\ \hline
Système de fichiers persistant & \pictoPleinAudit & \pictoPleinAudit & \pictoVideAudit \\ \hline
Isolation réseau fine & \pictoPartielAudit & \pictoPleinAudit & \pictoPartielAudit \\ \hline
Tâches longues & \pictoPleinAudit & \pictoPleinAudit & \pictoPartielAudit \\ \hline
Démarrage à froid & \pictoPleinAudit & \pictoPleinAudit & \pictoVideAudit \\
\hline
\end{tabular}
\end{table}

\subsection{Alternatives évaluées puis écartées, couche par couche}
\label{ann:alternatives-couches}

Cet inventaire complète le \cref{tab:briques-structurantes} de la section suivante : pour chacune des six couches, il
nomme les alternatives réellement évaluées avant la brique retenue. Il est publié ici parce qu'il
relève du matériel de reproduction --- il permet de refaire le choix --- et non de
l'argumentation, portée par le corps.

{\small\noindent
\emph{N1a} --- deux autres grands fournisseurs de cloud, un fournisseur européen généraliste,
deux hébergeurs à forfait ; le réseau par défaut ou partagé, le connecteur d'accès sans serveur,
le mandataire filtrant en sortie, le point de terminaison privé ; orchestrateur de conteneurs
managé, machines virtuelles, fonctions à la demande ; répartiteur régional, filtrage
auto-hébergé. \emph{N1b} --- identité unique partagée et droits au niveau du projet, clé de
compte de service exportée, agent et coffre auto-hébergés, chiffrement géré par le fournisseur,
module matériel de sécurité, fournisseur d'identité externe, facteur résistant à l'hameçonnage.
\emph{N1c} --- image mono-étage ou minimale, registre public ou auto-hébergé, cadre synchrone
classique, application monopage, migrations à la main. \emph{N1d} --- trois autres familles de
bases, les quatre autres modèles de supervision du \cref{tab:supervision-cinq}, deux ordonnanceurs
externes, disque local, agent de collecte tiers. \emph{N1e} --- format de règles neutre avec
moteur de conversion, moteur du fournisseur, modèle de chaîne d'attaque, modèle généraliste,
pondération de termes sans apprentissage de domaine, modèle génératif, inférence managée.
\emph{N1f} --- service d'intégration du fournisseur, moteur auto-hébergé, autre analyseur de
secrets, analyseur commercial, analyseurs d'images concurrents.\par}

\subsection{Sources et réserves de la comparaison d'hébergement}
\label{ann:etudes-comparatives-plateforme}

Les paliers gratuits (critères 1 et 3 du \cref{tab:plateforme-10-criteres}) et les implantations
régionales (critère 4) reposent sur la documentation publique des fournisseurs, consultée le
27/08/2026 : paliers gratuits et régions de Google
Cloud~\cite{gcp_free_tier,gcp_regions_zones}, régions d'Amazon Web Services~\cite{aws_regions},
implantations de Hetzner~\cite{hetzner_locations} et régions de Scaleway~\cite{scaleway_regions}.
Les colonnes \og{}Azure\fg{} et \og{}OVH\fg{} n'ont pas pu être sourcées à cette date, les pages de
documentation correspondantes étant dynamiques et inaccessibles depuis l'environnement de travail ;
leurs valeurs relèvent donc de la connaissance générale de l'auteur et non d'une vérification
documentaire.

Les deux réserves méthodologiques annoncées au \cref{sec:plateforme-criteres} trouvent ici leur
justification documentaire : aucune mesure de latence depuis Nouakchott n'ayant été conduite, le
critère~4 est apprécié à partir des seules implantations régionales publiées ci-dessus ; et les
grilles tarifaires n'ayant pas été relevées à une date unique ni pour un dimensionnement identique,
le critère~2 est apprécié à partir des modèles de facturation, non de prix mesurés.

\section{Inventaire technologique détaillé des trente-trois briques}
\label{ann:fiche-versions}

Cette annexe porte l'inventaire complet dont le \cref{sec:trente-trois-briques} ne publie que la
synthèse par couche. Elle se compose de deux tableaux qui décrivent \textbf{les mêmes
trente-trois briques dans le même ordre} et se lisent ligne pour ligne : le premier donne, pour
chacune, le critère qui a décidé de son adoption et le prix payé pour elle ; le second, la version
ou le paramètre retenu et le mode de verrouillage de cette version. Les deux colonnes décisives
sont respectivement le prix payé --- un avantage sans contrepartie nommée n'est pas un arbitrage
--- et le mode de verrouillage --- une version connue ne dit rien de la reproductibilité si rien
n'empêche la valeur de changer seule.

\subsection{Critère décisif et prix payé, brique par brique}
\label{ann:briques-detail}

Le rôle de chaque brique est porté par son libellé, le produit étant nommé entre parenthèses à sa
première occurrence.

\begingroup
% NB : la taille est imposee a 8 pt par \AtBeginEnvironment{longtable}{\scriptsize}
% (config/formatting.tex) ; toute declaration de taille posee ici serait sans effet.
\setlength{\tabcolsep}{2pt}
\renewcommand{\arraystretch}{0.92}
\setlength{\LTpre}{2pt}
\setlength{\LTpost}{2pt}
\begin{longtable}{|P{3.50cm}|P{5.10cm}|P{6.90cm}|}
\caption{Les trente-trois briques structurantes, par couche : critère décisif et prix payé}
\label{tab:briques-structurantes}\\
\hline
\textbf{Brique et rôle} & \textbf{Critère décisif} & \textbf{Prix payé} \\
\hline
\endfirsthead
\hline
\textbf{Brique et rôle} & \textbf{Critère décisif} & \textbf{Prix payé} \\
\hline
\endhead
\hline
\endfoot
\multicolumn{3}{|l|}{\textbf{N1a --- Socle et infrastructure : huit briques}} \\ \hline
Fournisseur de cloud public unique (Google Cloud) : support de l'ensemble du socle, dans une
région européenne unique &
Seul à réunir dans un même service le stockage, le langage de détection et la recherche
vectorielle (\cref{sec:plateforme-criteres}) &
Réversibilité perdue, coût direct supérieur à un forfait, aucune région en Afrique de
l'Ouest \\ \hline

Infrastructure décrite en code (Terraform, 16~modules) : provisionnement, reproductibilité,
détection de dérive &
Maturité de l'écosystème de modules et portabilité de la description &
Langage déclaratif : toute logique conditionnelle passe par des compteurs, moins lisibles qu'un
langage général. L'état est un actif à protéger que le code ne décrit pas, et un plan ciblé peut
inclure des remplacements destructifs sans rapport \\ \hline

Exécution de conteneurs sans serveur (Cloud Run, D01) : support des cinq services et des trois
tâches &
Aucun système d'exploitation à durcir ni cluster à exploiter --- décisif sous contrainte
mono-opérateur &
Démarrage à froid mesuré : environ deux secondes sur l'API, 94~s au centile~99 sur le service
d'encodage, ce qui a imposé de recalibrer le budget de la sonde. Aucune charge à état, aucune
tâche longue \\ \hline

Réseau privé virtuel en mode personnalisé : un sous-réseau par locataire, accès privé aux
données, refus par défaut dans les deux sens &
Contrôle explicite des plages et refus par défaut \textbf{en entrée et en sortie}, chaque
autorisation étant nominative : aucune règle d'autorisation ne comporte de plage source, ce qui
serait de la confiance de zone &
Trois sous-réseaux et douze règles là où l'ancien plan en déclarait deux et quatre : les refus
croisés croissent comme le carré du nombre de locataires, ce qui borne la conception à trois ou
quatre. Et le sélecteur demeure l'\textbf{étiquette réseau}, la plateforme n'exposant pas le
compte de service au niveau des interfaces de sortie : la frontière est opposable à une charge de
travail compromise, \textbf{non à un opérateur capable de déployer} \\ \hline

Sortie réseau directe : interface réseau attachée à chaque instance dans le sous-réseau de son
locataire, en remplacement du connecteur d'accès sans serveur (D15, bascule service par service
du 12 au 15/08/2026) &
C'est ce qui rend le pare-feu applicable aux charges de travail : sans adresse dans un
sous-réseau, aucune règle ne peut les désigner, ni comme source ni comme cible, et un refus par
défaut en sortie resterait une déclaration sans effet &
Chaque instance active consomme une adresse du sous-réseau, là où le connecteur en immobilisait
seize quel que fût le trafic : une montée en charge imprévue peut épuiser un \texttt{/26}. Le
rattachement de l'interface intervient au démarrage de l'instance, donc dans le chemin du
démarrage à froid. Enfin le socle \textbf{perd un objet observable} : les interfaces sont
éphémères et ne se donnent à voir que par les journaux de flux \\ \hline

Accès privé aux services managés : connexion à la base sans adresse publique &
Élimine structurellement l'exposition de la base &
Plage allouée automatiquement par le fournisseur, non maîtrisée dans la description ; le module
de base n'est pas réinstanciable en l'état, le pairage étant câblé à l'intérieur \\ \hline

Répartiteur de charge applicatif global : point d'entrée public unique, terminaison du
chiffrement, routage par nom d'hôte vers quatre services &
Point d'entrée unique = point de contrôle unique ; certificats gérés et renouvelés par le
fournisseur &
La gestion avancée du trafic --- pondération, mise en miroir, politiques de réessai --- n'est pas
disponible dans la variante retenue. L'entrée des services étant restreinte à ce seul chemin, la
garantie repose sur un contrôle unique, sans défense en profondeur \\ \hline

Filtrage applicatif managé au bord (Cloud Armor) : refus des charges d'attaque connues,
limitation de débit, restriction géographique &
Aucun composant à exploiter ; règles maintenues par le fournisseur ; application uniforme aux
quatre points d'entrée, sans exception &
Cinq familles de règles retenues sur onze disponibles, à la sensibilité la plus basse :
arbitrage explicite entre faux positifs et évasions --- quatre motifs sur cinq bloqués au test
réel du 19/08/2026. Aucune protection adaptative : le seuil anti-saturation est fixe \\ \hline

\multicolumn{3}{|l|}{\textbf{N1b --- Identité, secrets et chiffrement : six briques}} \\ \hline
Gestion des identités et des accès du fournisseur : sept identités de service, droits au moindre
privilège, aucun rôle primitif &
Une identité par charge de travail : c'est ce qui remplace les pare-feux internes (PR2) &
Plusieurs liaisons demeurent à l'échelle du projet là où une portée par ressource serait possible ;
la granularité par instance n'existe pas pour la base ; et aucune politique d'organisation ne
rend la règle préventive (É7) --- elle est tenue par discipline, vérifiée par inventaire \\ \hline

Fédération d'identité pour la chaîne de livraison (D08) : authentification de la chaîne sans clé
permanente, sous condition de dépôt et de branche &
Supprime la classe de risque au lieu de la gérer : il n'y a pas de clé à faire fuir &
Dépendance à la forge logicielle comme émetteur d'assertion : sa compromission devient un chemin
d'accès. La restriction de branche tient par une condition d'attribut, non par la liaison
elle-même \\ \hline

Coffre de secrets managé : neuf secrets, un par usage, injectés par référence dans les services &
Chiffrement par clé gérée et liaison d'accès \textbf{par secret}, jamais au niveau du projet &
Dépendance réseau à l'exécution pour un secret, là où le montage natif était disponible ; un
secret subsiste hors clé gérée, vestige d'une interface antérieure dont le retrait est
planifié \\ \hline

Service de gestion de clés, chiffrement par clés gérées : entrepôt, secrets et espace média ;
rotation à 90~jours &
Contrôle de la rotation et de la révocation, exigé par la nature des données de l'application
hébergée &
Niveau de protection logiciel et non matériel ; le chiffrement de l'entrepôt n'est pas
rétroactif ; il est structurellement inapplicable à la base, le champ étant immuable après
création \\ \hline

Authentification applicative par jeton signé et rôles (D10) : contrôle d'accès à l'API de
supervision, trois rôles, vérification route par route &
Aucune dépendance externe ; le rôle devient visible dans la signature de chaque point d'entrée,
donc auditable en lecture &
Aucun mécanisme de révocation : un jeton reste valide jusqu'à son expiration après une
déconnexion ou un changement de rôle, et l'état actif d'un compte n'est vérifié qu'à
l'émission \\ \hline

Second facteur par code temporel à usage unique : renforcement de l'accès au tableau de bord ;
secret chiffré au repos &
Standard ouvert, vérifiable hors ligne, sans dépendance ni coût par usage &
Pas de résistance à l'hameçonnage ; aucun code de secours, donc la perte du terminal exige une
intervention en base ; le secret est protégé par un chiffrement applicatif dont la clé unique
sans version interdit toute rotation \\ \hline

\multicolumn{3}{|l|}{\textbf{N1c --- Exécution, conteneurisation et applications : six briques}} \\ \hline
Conteneurisation multi-étage, utilisateur non privilégié, base épinglée par empreinte : format de
livraison de tous les services &
Surface réduite et image finale débarrassée des outils de construction, ce qui a permis
d'éliminer une vulnérabilité critique en retirant le gestionnaire de paquets &
Construction plus longue et fichiers de description plus complexes ; la mise à jour de la base
épinglée redevient un geste manuel à chaque publication \\ \hline

Registre d'images privé du fournisseur (D07) : stockage des artefacts déployés &
Un registre unique : deux registres constituent deux surfaces de chaîne d'approvisionnement et
une ambiguïté sur la source de vérité &
Aucune immuabilité d'étiquette déclarée, aucune politique de nettoyage, et l'analyse de
vulnérabilité reste externe à la plateforme --- l'image y est publiée avant d'être
analysée \\ \hline

Cadre applicatif à validation typée pour l'API : vingt et un points d'entrée, unique accès aux
données de sécurité &
Validation déclarative en entrée \textbf{et} contrat de sortie forcé : le contrôle d'accès
devient visible dans la signature de chaque point d'entrée &
Aucun bénéfice de l'asynchronisme, les pilotes retenus étant synchrones. Aucune version portée
par les points d'entrée ; la documentation d'interface demeure publique \\ \hline

Cadre web à rendu serveur pour le tableau de bord (D14) : douze vues, sept routes serveur &
Le jeton d'authentification reste côté serveur et n'est jamais lisible par un script client
--- propriété de sécurité, pas de confort &
Rendu dynamique à chaque requête, aucune mise en cache statique ; démarrage à froid au-delà de
dix secondes, qui impose d'y maintenir une instance active ; couverture de tests de
3,00~\% au 28/08/2026 \\ \hline

Accès à la base par connecteur applicatif chiffré : liaison entre l'API et la base &
Tunnel chiffré mutuellement authentifié vers l'adresse privée : c'est le plus strict des trois
modes disponibles &
Dépendance à une bibliothèque cliente supplémentaire et à un pilote purement interprété, donc
synchrone --- c'est ce pilote qui interdit l'asynchronisme de la ligne précédente \\ \hline

Outil de migration de schéma versionné : évolution du schéma relationnel de la plateforme &
Le schéma devient reconstructible depuis le dépôt --- réponse directe au constat bloquant B5 de
l'audit &
Aucune automatisation dans la chaîne : les migrations de la plateforme sont déclenchées à la
main, là où l'application hébergée les exécute par une tâche avant déploiement \\ \hline

\multicolumn{3}{|l|}{\textbf{N1d --- Données et entrepôt de supervision : cinq briques}} \\ \hline
Base relationnelle managée, haute disponibilité régionale (D03, D13) : comptes, rôles, journal
d'audit applicatif, données de l'application hébergée &
Restauration à un instant donné et bascule régionale sans exploitation à charge --- décisif
sous contrainte mono-opérateur, et mesuré à 32~min~45~s pour zéro perte le 03/08/2026 &
Instance de gamme d'entrée facturée au tarif haute disponibilité ; chiffrement par clés gérées
structurellement impossible ; \textbf{une seule instance partagée entre les locataires} \\ \hline

Entrepôt de données généraliste comme système de supervision (BigQuery, D04) : dix tables ---
journaux normalisés, détections, enrichissement, vulnérabilités, verdicts &
Seule option offrant stockage, langage de détection \textbf{et} recherche vectorielle dans le
même service, sans base spécialisée (D06) &
Ni règles, ni interface, ni corrélation, ni support fournis : trois chantiers à construire.
Plafond de règles maintenables par une personne. Poste de coût dominant, avec un minimum
forfaitaire par table balayée. Aucun index vectoriel déployé : chaque recherche parcourt tout le
référentiel \\ \hline

Requêtes planifiées de l'entrepôt comme ordonnanceur : douze exécutions, cinq normalisations et
sept règles de détection &
Aucun composant supplémentaire : l'ordonnancement vit là où vivent les données (PR1), et le
moteur d'orchestration a été retiré après soixante exécutions en échec consécutives &
Cadence plancher de cinq minutes imposée par le service ; aucune dépendance entre exécutions,
donc aucune garantie d'ordre ; les seuils sont figés dans les requêtes \\ \hline

Stockage objet avec prévention d'accès public et clés gérées : espace média de l'application
hébergée &
L'application d'origine écrivait sur le disque local, ce qui interdisait toute mise à l'échelle
et toute reprise : c'était le constat bloquant B6 &
Aucun versionnement d'objet, aucune règle de cycle de vie : une suppression accidentelle est
irréversible \\ \hline

Collecte de journaux et exports vers l'entrepôt : quatre exports déclarés en code, filtrage à la
source &
Filtrage \textbf{au moment de la collecte} et non après stockage : c'est le levier de coût
principal (besoin BNF3) &
Trois exports sur quatre filtrent réellement ; celui du répartiteur ingère la totalité du trafic
et reste le poste de coût dominant. La liste des services collectés étant explicite, un service
oublié disparaît de la supervision sans erreur visible. Les journaux d'accès aux données restent
dans l'espace par défaut, à conservation modifiable \\ \hline

\multicolumn{3}{|l|}{\textbf{N1e --- Détection et enrichissement sémantique : quatre briques}} \\ \hline
Règles de détection versionnées, écrites dans le dialecte de l'entrepôt : sept règles R1 à R7,
fenêtre glissante, déduplication, étiquetage par technique &
Aucun moteur de conversion ne cible ce dialecte : la traduction manuelle évite une dépendance
de conversion pour sept règles (verrou V2) &
Aucune portabilité vers un autre système de supervision, aucun accès au corpus communautaire,
et un plafond de couverture lié au volume maintenable par une personne --- six techniques et
cinq tactiques sur quinze \\ \hline

Référentiel de techniques d'attaque (MITRE ATT\&CK) : vocabulaire commun conception-détection,
grille de couverture, cible de l'enrichissement &
Seul référentiel qui fournit à la fois un vocabulaire de détection et un catalogue textuel
encodable &
Le dénominateur est la matrice intégrale, qui couvre des systèmes absents du périmètre : tout
taux rapporté à ce dénominateur est structurellement bas, ce qui conduit le mémoire à publier
l'inventaire nominatif des techniques plutôt qu'un pourcentage \\ \hline

Modèle de représentation de phrases spécialisé en cybersécurité (D05) : encodage des détections
et des techniques en vecteurs comparables &
Représentation par phrase entraînée sur le vocabulaire du domaine ; \textbf{et surtout}, un
modèle qui produit un vecteur n'exécute aucune instruction, là où soumettre des journaux à un
modèle génératif transmettrait une entrée contrôlée par l'attaquant à un composant qui suit des
instructions &
Aucune explication en langue naturelle ; poids d'environ 440~mégaoctets embarqués dans l'image ;
consommation mémoire au pic de l'ordre de 900~mégaoctets ; qualité non établie sur des alertes
courtes et bruitées (verrou V3) \\ \hline

Moteur d'inférence embarqué, exécution sur processeur généraliste : service d'encodage interne,
sans sortie publique et sans droit sur l'entrepôt &
Exécution locale et déterministe, sans appel externe ni coût par appel, et sans transmission de
données de sécurité à un tiers &
Démarrage à froid de 94~s au pic, qui a imposé de recalibrer le budget de la sonde de
disponibilité ; l'export quantifié a été testé puis rejeté sur mesure, donc les poids restent
en précision flottante simple \\ \hline

\multicolumn{3}{|l|}{\textbf{N1f --- Chaîne de livraison et chaîne de provisionnement : quatre
briques}} \\ \hline
Forge logicielle et son moteur d'exécution (GitHub Actions) : porte la chaîne de livraison
applicative (F3) et la chaîne de provisionnement de l'infrastructure (F7) &
Intégration native avec la fédération d'identité, qui supprime la clé permanente &
Les actions réutilisées sont référencées par étiquette mobile : une redirection d'étiquette
ferait exécuter un code différent sans qu'une ligne du dépôt change \\ \hline

Détection de secrets versionnés (Gitleaks, D09) : première porte bloquante de la chaîne &
Coût quasi nul, gain élevé, et \textbf{placement en tête de chaîne} : un secret est compromis
dès sa publication, le détecter tôt déclenche la rotation pendant qu'elle sert encore &
Jeu de règles par défaut, non spécialisé au domaine ; résultats disponibles dans les journaux
de la chaîne uniquement, sans format d'échange normalisé \\ \hline

Analyse statique du code (Semgrep) : deuxième porte bloquante &
Exécutable localement, jeu de règles ouvert, exclusions justifiables ligne à ligne plutôt que
globalement &
Le jeu de règles distant évolue : le verdict pouvait varier à version d'outil identique, ce qui
a imposé un verrouillage de version. Sept exclusions nominatives à maintenir \\ \hline

Analyse de vulnérabilité des images et de la configuration (Trivy) : troisième porte bloquante,
sur les deux chaînes, avec restitution du rapport dans la supervision (F6) &
\textbf{Un seul outil couvre les images et la description d'infrastructure} : un composant en
moins, application directe de PR1 &
Aucune génération de nomenclature logicielle ni de signature d'artefact, alors que l'outil en
est capable : deux familles de contrôle restent découvertes (décision D07). Les vulnérabilités
sans correctif amont sont ignorées \\
\end{longtable}
\endgroup

\subsection{Version retenue et mode de verrouillage}
\label{ann:versions-detail}

Le \cref{tab:fiche-versions} donne, brique par brique et dans l'ordre du tableau précédent, la
version ou le paramètre retenu et le \textbf{mode de verrouillage}. Les valeurs ont été relevées le \textbf{28/08/2026}, à deux exceptions signalées en
cellule : empreintes des images de base (19/08/2026) et version de l'analyseur statique
(24/08/2026).

\emph{Verrouillé} : un
fichier de verrouillage suivi dans le dépôt fige la valeur. \emph{Épinglé} : la valeur est écrite
dans le code, lisible et revue, mais rien n'en atteste l'intégrité. \emph{Étiquette mobile} : la
valeur change sans qu'une ligne du dépôt soit modifiée. \textbf{Une seule ligne est dépourvue de
version relevable}, en fin de tableau : la configuration de l'environnement de recette --- chemins
protégés, services collectés, adresses d'administration exemptées --- est portée par un fichier de
valeurs exclu du dépôt, qui contient des décisions de sécurité et aucun secret. Son versionnement
est la mesure de reproductibilité la plus rentable qu'ait identifiée le relevé.

\begingroup
\scriptsize
\setlength{\tabcolsep}{3pt}
\renewcommand{\arraystretch}{1.0}
\setlength{\LTpre}{4pt}
\setlength{\LTpost}{4pt}
\begin{longtable}{|P{3.4cm}|P{7.5cm}|P{4.4cm}|}
\caption[Fiche de version des briques structurantes]{Fiche de version des trente-trois briques structurantes et mode de verrouillage}
\label{tab:fiche-versions}\\
\hline
\textbf{Brique} & \textbf{Version ou paramètre retenu} & \textbf{Verrouillage} \\
\hline
\endfirsthead
\hline
\textbf{Brique} & \textbf{Version ou paramètre retenu} & \textbf{Verrouillage} \\
\hline
\endhead
\hline
\endfoot
\multicolumn{3}{|l|}{\textbf{N1a --- Socle et infrastructure}} \\ \hline
Fournisseur de cloud public unique &
Projet unique, région \texttt{europe-west1} &
Épinglé (variables de la racine d'environnement) \\ \hline
Infrastructure décrite en code &
Outil \texttt{>= 1.7} (borne basse seule) ; fournisseurs de ressources 5.45.2, aléatoire 3.9.0,
forge 6.13.0 &
Fournisseurs verrouillés par fichier versionné ; l'outil n'a pas de borne haute \\ \hline
Exécution de conteneurs sans serveur &
Deuxième génération exclusivement ; sortie intégrale vers le sous-réseau du locataire, sortie
managée pour les deux charges détachées ; entrée limitée au répartiteur &
Épinglé (type de ressource et paramètres codés dans le module) \\ \hline
Réseau privé virtuel personnalisé &
Trois sous-réseaux en /26, deux occupés ; accès privé aux services du fournisseur ; douze règles de
pare-feu et refus par défaut en priorité 65534 en entrée comme en sortie
(\cref{subsec:reseau-cible}) ; journaux de flux échantillonnés à 0,1, métadonnées exclues &
Épinglé \\ \hline
Sortie réseau directe &
Une interface par instance dans le sous-réseau de son locataire ; sortie \texttt{ALL\_TRAFFIC} ;
une étiquette réseau par compte de service &
Épinglé (sous-réseau et étiquette codés dans le module). La solidarité entre l'étiquette et le
compte de service tient à leur déclaration dans le même bloc, donc à une convention d'écriture et
non à un contrôle ; la bascule vers le sélecteur par compte de service est prévue par une variable
du module \\ \hline
Accès privé aux services managés &
Plage réservée de longueur /16 &
Partiellement épinglé : préfixe déclaré, adresse de début choisie par le fournisseur \\ \hline
Répartiteur de charge global &
Variante classique ; profil de chiffrement moderne, version minimale de transport 1.2 ; trois
certificats gérés &
Épinglé \\ \hline
Filtrage applicatif managé &
Neuf règles ; cinq jeux de signatures en version \texttt{v33-stable}, sensibilité 1 ; seuils de
10 et de 1000 requêtes par 60~s &
Épinglé quant au nom de version ; contenu des signatures maintenu par l'éditeur, non
lisible \\ \hline

\multicolumn{3}{|l|}{\textbf{N1b --- Identité, secrets et chiffrement}} \\ \hline
Gestion des identités et des accès &
Sept identités de service ; aucun rôle primitif ; aucune clé exportée &
Épinglé (liaisons déclarées en code) \\ \hline
Fédération d'identité de la chaîne &
Condition d'attribut portant sur le dépôt \emph{et} sur la branche \texttt{refs/heads/main} &
Épinglé \\ \hline
Coffre de secrets managé &
Neuf secrets, huit sous clé gérée &
Épinglé ; versions de secret référencées par \texttt{latest}, donc mobiles \\ \hline
Service de gestion de clés &
Deux clés ; rotation à 90 jours ; niveau de protection logiciel &
Épinglé \\ \hline
Authentification par jeton signé et rôles &
Algorithme de signature figé en constante ; durée de vie 60~min ; jeton intermédiaire 5~min sans
rôle &
Épinglé dans le code applicatif \\ \hline
Second facteur par code temporel &
Standard RFC~6238 ; secret de 160~bits ; fenêtre de dérive d'un pas &
Épinglé \\ \hline

\multicolumn{3}{|l|}{\textbf{N1c --- Exécution, conteneurisation et applications}} \\ \hline
Conteneurisation multi-étage, utilisateur non privilégié &
Bases épinglées par empreinte : \texttt{sha256:a116514e\ldots} (interpréteur Python 3.12),
\texttt{sha256:fb4cd12c\ldots} (exécution Node 20), \texttt{sha256:2c941e86\ldots} (service
d'encodage) --- relevé le 19/08/2026 &
Verrouillé par empreinte dans les fichiers de construction versionnés \\ \hline
Registre d'images privé &
Format conteneur, région \texttt{europe-west1} ; aucune étiquette immuable déclarée &
Épinglé quant au registre ; artefacts de l'application hébergée référencés par étiquette
mobile \\ \hline
Cadre applicatif de l'API &
Couche d'accès aux données 2.0.36 ; dépendances figées &
Verrouillé par fichier de dépendances versionné \\ \hline
Cadre web du tableau de bord &
14.2.35 &
Verrouillé : l'installation se fait à partir du fichier de verrouillage versionné \\ \hline
Connecteur applicatif chiffré vers la base &
Type d'adresse privé imposé ; réserve de 5 connexions, débordement de 2, recyclage à 1800~s &
Épinglé \\ \hline
Outil de migration de schéma &
Trois révisions linéaires, de \texttt{001} à \texttt{003} &
Verrouillé par fichier versionné : les révisions vivent dans le dépôt \\ \hline

\multicolumn{3}{|l|}{\textbf{N1d --- Données et entrepôt de supervision}} \\ \hline
Base relationnelle managée &
PostgreSQL 15 ; classe de machine à cœur partagé ; haute disponibilité régionale ; retour à un
instant donné sur 7 jours ; 100 connexions &
Épinglé \\ \hline
Entrepôt de supervision &
Un jeu de données, dix tables ; partitionnement par jour ; expiration à 90 jours sur les deux
tables de journaux &
Épinglé \\ \hline
Requêtes planifiées comme ordonnanceur &
Douze exécutions à cadence de cinq minutes ; fenêtre glissante de quinze minutes &
Épinglé \\ \hline
Stockage objet &
Accès uniforme au niveau du dépôt, prévention d'accès public appliquée, clé gérée régionale &
Épinglé \\ \hline
Collecte de journaux et exports &
Quatre exports ; identité d'écriture unique par export ; tables partitionnées &
Épinglé \\ \hline

\multicolumn{3}{|l|}{\textbf{N1e --- Détection et enrichissement sémantique}} \\ \hline
Règles de détection versionnées &
Sept règles, R1 à R7, définies dans un module unique &
Verrouillé par fichier versionné \\ \hline
Référentiel de techniques d'attaque &
Version 19 (28/04/2026), quinze tactiques &
Consigné dans le mémoire ; le chargement du catalogue ne fige aucune version dans le code \\ \hline
Modèle de représentation de phrases &
Export en précision flottante simple, étiquette \texttt{v1.0}, 768 dimensions, seuil de
similarité 0,60 &
Épinglé et vérifié à l'exécution : la tâche refuse de s'exécuter hors développement sur
l'étiquette d'un substitut \\ \hline
Moteur d'inférence embarqué &
Exécution sur processeur généraliste ; deux gibioctets de mémoire ; budget de démarrage de
150~s &
Épinglé \\ \hline

\multicolumn{3}{|l|}{\textbf{N1f --- Chaîne de livraison et chaîne de provisionnement}} \\ \hline
Forge logicielle et moteur d'exécution &
Actions de récupération du dépôt, d'authentification et de déploiement en \texttt{v4} et
\texttt{v2} &
Étiquette mobile --- dette reconnue, écartée nominativement de l'analyse statique, justification
écrite \\ \hline
Détection de secrets versionnés &
Action de forge \texttt{v2} ; jeu de règles par défaut &
Étiquette mobile \\ \hline
Analyse statique du code &
1.173.0 ; jeu de règles par défaut de l'éditeur ; sept exclusions nominatives --- relevé le
24/08/2026 &
Épinglé pour l'outil ; le jeu de règles est distant et mobile \\ \hline
Analyse de vulnérabilité des images et de la configuration &
Action de forge \texttt{v0.36.0} ; blocage sur criticité maximale corrigeable &
Épinglé \\ \hline

\multicolumn{3}{|l|}{\textbf{Paramétrage transverse des briques ci-dessus}} \\ \hline
Configuration de l'environnement de recette (chemins protégés, services collectés, adresses
d'administration) &
--- &
Non versionné : fichier de valeurs exclu du dépôt bien qu'il ne contienne aucun secret \\
\end{longtable}
\endgroup

\medskip

\begin{alertbox}[boxsep=1pt, left=4pt, right=4pt, top=3pt, bottom=3pt]
\textbf{Ce que le verrouillage des versions ne verrouille pas.} Épingler la version d'un outil
d'analyse ne fige pas le \textbf{jeu de règles distant} qu'il télécharge à l'exécution, comme
l'établit le décompte de règles chargées les 20 et 24/08/2026 à version d'outil identique
(\cref{sec:chaine-livraison}) : la chaîne est reproductible quant aux \emph{outils} qu'elle
exécute, non quant aux \emph{critères} qu'ils appliquent. Le contrôle compensatoire est la
conservation des journaux d'exécution, qui portent le nombre de règles évaluées ; le seuil de
bascule serait l'hébergement du jeu de règles dans le dépôt. Le raisonnement vaut aussi pour les
signatures du filtrage applicatif et les versions de secret référencées par étiquette mobile.
\end{alertbox}
